\section*{Appendix A: Responsible NLP Research Checklist}

The Responsible NLP Research Checklist below is filled per ARR / EMNLP submission guidelines.

\paragraph{B1. Scientific artifacts.}
\begin{itemize}\setlength\itemsep{1pt}
\item \textbf{Cited and used:} HotpotQA \citep{yang2018hotpotqa} under its public Apache-2.0 license; \texttt{gpt-4.1-mini} accessed via an OpenAI-compatible API endpoint per its commercial terms of service.
\item \textbf{License/intent compatibility:} HotpotQA license permits research use. The model API terms permit benchmarking and reporting of aggregate metrics, which is the entirety of our model use.
\item \textbf{Intended-use compatibility:} our use is non-commercial multi-hop reading-comprehension benchmarking, consistent with HotpotQA's stated scope.
\item \textbf{Released artifacts:} runtime code (\texttt{runner.py}, \texttt{methods.py}, \texttt{contracts.py}), seed slice, validation scripts, and per-run logs are released in the anonymous code/data supplement under MIT license. Use is restricted to research; no derived models are released.
\end{itemize}

\paragraph{B2. Computational experiments.}
\begin{itemize}\setlength\itemsep{1pt}
\item \textbf{Model size / parameters:} we do not train or fine-tune any model; the only model used is the API-served \texttt{gpt-4.1-mini} (size not disclosed by provider).
\item \textbf{Compute budget:} entire experiment requires $\le 10^4$ API calls per method per chain-200 batch (i.e., 4 hops $\times$ 200 examples $\times$ 3 methods $\approx 2{,}400$ contribution calls + 600 answer calls). \textbf{Wall-clock estimate:} approximately 1--2 hours per chain-200 batch under the parallel runner (subject to provider rate-limit). \textbf{USD cost order-of-magnitude:} approximately \$0.50--\$1.00 per chain-200 batch with the canonical \texttt{gpt-4.1-mini} backbone (3 methods); the historical glm-4-flash backbone runs were closer to \$0.10 per batch. The full-validation pass ($n{=}7405$) is approximately $35{\times}$ a chain-200 batch on the same backbone. No GPU compute on our side.
\item \textbf{Hyperparameters:} routing weights for the Stage-1 TCPB instance ($0.55, 0.20, 0.10$, etc.), accept margin ($\delta_a{=}0.02$), update step sizes ($\Delta^{+}{=}{+}0.06$, $\Delta^{-}{=}{-}0.10$), momentum ($\rho{=}0.5$), F1 threshold ($F_\textsc{th}{=}0.5$), max handoff ($H_{\max}{=}4$). The Stage-2 mechanism constants ($\nu{=}0.2$, $k_{\max}{=}3$, $d_{\max}{=}3$, $n_{\max}{=}12$) are listed in Algorithm~\ref{alg:edo-stage2-loop}; all are hand-set priors, not searched.
\item \textbf{Random seeds:} all chain-200 results in the main paper tables use \texttt{seed=42} for both example slicing (deterministic first-200 of the HotpotQA validation split) and any internal stochasticity (\texttt{Python random.seed(42)}); the exact 200 question IDs are listed in the anonymous code/data supplement.
\item \textbf{Reported summary statistics:} EM, F1, MHC, PAR, and per-sample API token counts (reported in the main paper tables). The local \texttt{Phi-4-mini} diagnostic reports a paired-bootstrap CI; the strong-backbone Stage-1 table still lacks paired CIs and is flagged in \S Limitations as statistically incomplete.
\item \textbf{Existing implementations used:} we use the standard HotpotQA EM/F1 evaluator (re-implemented in our codebase to match the original word-tokenisation rules); no other off-the-shelf evaluators are used.
\end{itemize}

\paragraph{B3. Human subjects / annotators.} Not applicable. No human annotation was collected for this work, and no human evaluation of model outputs was conducted. The only human-curated data used is the public HotpotQA dataset (and its Wikipedia-sourced gold answers / supporting facts).

\paragraph{B4. Use of AI assistants.} The first author used a code-completion AI assistant during the development of \texttt{runner.py}, \texttt{methods.py}, and prompt templates. All AI-suggested code was reviewed line-by-line by the author before commit, and the final algorithmic design (Algorithm~\ref{alg:edo-stage2-loop}, Prototype Scope Box) is the author's own contribution. Manuscript prose, claims, citations, and reported numbers were authored and checked by the authors; any AI assistance for wording was limited to grammar/style phrasing and did not introduce scientific content.

\paragraph{B5. Use of data containing personally identifiable information (PII).} Not applicable. The HotpotQA benchmark used in this work is sourced from publicly released Wikipedia article excerpts and the question--answer pairs constructed over them; the dataset contains no personally identifiable information about article subjects beyond what is already publicly available on Wikipedia, and no additional PII was collected, generated, or processed at any stage of this work.

\section*{Appendix B: Provider Integrity Event --- Engineering Response}
\label{sec:appendix-b}

This appendix documents the engineering response to the provider-side data-integrity event referenced in \S Limitations item~(5). It is placed in the appendix per the EMNLP/ACL convention that Limitations should describe the scope and consequence of issues, while the corresponding engineering implementation belongs in supplementary material.

\paragraph{Event description.} During the full-validation pass on the strong backbone, two of the three method runs were silently affected by an upstream provider issue in which the requested model identifier was routed to a different (newer) underlying model for a subset of API calls; the third method's full-validation result was not affected.

\paragraph{Runtime guard.} The runtime client was extended with a fail-fast integrity check that compares the model identifier echoed back in each API response against the requested identifier and raises an explicit error on mismatch. The check operates per-call, before any downstream consumption of the response content.

\paragraph{Endpoint switch.} The upstream API endpoint serving the strong backbone was switched to a different OpenAI-compatible provider channel with the integrity guard enabled in pre-send mode: a known-fake model identifier raises immediately without issuing a request, demonstrating the guard's pre-flight effectiveness.

\paragraph{Affected results disposition.} The two affected method runs are queued for re-execution under the new guard and the new endpoint. The third (unaffected) full-validation result reported in the main paper is preserved as-is, since invalidating it would discard reproducible evidence without obtaining new information. Cross-endpoint comparability between the original main-paper numbers and the deferred re-runs is the corresponding scope-of-claims caveat recorded in \S Limitations item~(5).

\section*{Appendix C: LLM Prompt Templates}
\label{sec:appendix-c}

For reproducibility, this appendix lists the four LLM-call templates referenced by Algorithm~\ref{alg:edo-stage2-loop}: \textsc{LLM\_Answer} (line 13), \textsc{LLM\_Decompose} (line 26), \textsc{Audit} (line 7, optional LLM-fallback variant), and \textsc{Evidence\_Extract} (line 33). All four are released verbatim in the anonymous code/data supplement; abridged versions are reproduced below.

\paragraph{C.1 \textsc{LLM\_Answer} (per-agent answer/forward template).}
A shared prefix establishes the global contract (local visibility only, plain-text output, HotpotQA short-answer style); per-role suffixes refine output for \texttt{decomposer}, \texttt{evidence\_seeker}, \texttt{verifier}, and \texttt{synthesizer} nodes. The synthesizer follows a two-stage internal protocol: (i) a brief "Key facts" recap reading all evidence lines from first to last, then (ii) deriving a single shortest answer span. The forward branch (when an agent forwards rather than answers) outputs exactly one short sentence adding one new fact or subgoal-relevant clue.

\paragraph{C.2 \textsc{LLM\_Decompose} (R1 \texttt{split} template).}
Takes one complex question and outputs at most $k_{\max}{=}3$ atomic, independently-answerable subquestions in strict JSON. Hard requirements: each subquestion is solvable in isolation (no inter-subquestion chain); each yields a concrete fact (no meta-questions); the original question is never reproduced as a subquestion; if no atomic decomposition is possible, the template returns \texttt{\{"subtasks": []\}} and the runtime downgrades to \textsc{DoSelf}. Each subquestion carries a \texttt{task\_type\_guess} hint $\in$ \{\texttt{factoid}, \texttt{comparison}, \texttt{aggregation}, \texttt{verification}\} for downstream routing. Two worked decomposition examples (a comparison and an aggregation) are included in the prompt body.

\paragraph{C.3 \textsc{Audit} (R2 audit decision template, LLM-fallback variant).}
The default audit is rule-based (no LLM call) per the Algorithm~\ref{alg:edo-stage2-loop} default. The optional LLM-fallback variant takes an upstream--downstream pair $(u, v)$, a task node $z$, and a candidate result $r$, and outputs a strict JSON object \texttt{\{decision, value\_gain, rework\_cost, rationale\}} with \texttt{decision} $\in$ \{\textsc{Accept}, \textsc{AcceptWithNote}, \textsc{RejectReroute}, \textsc{RejectResplit}\}. Decision rules (verbatim from the prompt): \textsc{Accept} when the candidate plausibly answers the subtask; \textsc{AcceptWithNote} when usable but incomplete or low-confidence; \textsc{RejectReroute} when wrong/empty AND another neighbour has a reasonable chance; \textsc{RejectResplit} when wrong AND the failure looks structural. \texttt{value\_gain}$\in[0,1]$ measures movement toward the root task; \texttt{rework\_cost}$\in[0,1]$ defaults to $1{-}\texttt{value\_gain}$.

\paragraph{C.4 \textsc{Evidence\_Extract} (R3 vector-belief update mapping).}
This is a deterministic rule-based mapping (no LLM call), not a prompt template. Given an audit event $e$ and the task signature $\phi(z)$, it projects $e$ onto the $7$-dimensional persona axes $[\text{solve}, \text{decompose}, \text{audit}, \text{integrate}, \text{explore}, \text{efficiency}, \text{reliability}]$ via a lookup-table from $e.\text{decision}$ and $e.\text{value\_gain}$ to per-axis EMA increments, weighted by $\phi(z)$'s axis prominence. The full mapping table appears in the anonymous executable specification supplement \S 4.

\section*{Appendix D: Preliminary Stage-2 Prototype 3-Shard Paired Result}
\label{sec:appendix-d}

For transparency we record a small-scale paired comparison of \texttt{fixed\_peer\_calibrated} (Stage-1) versus a degenerate-Stage-2 prototype (\texttt{edo\_stage2\_chain}, in which \textsc{Split} downgrades to a single inline \textsc{Outsource} step and \textsc{Audit} defaults to \textsc{Accept}) on three disjoint 200-sample shards of the HotpotQA validation split. Same backbone (\texttt{gpt-4.1-mini}), same chain topology, same handoff cap, paired questions per shard.

\begin{table}[h]
\centering
\small
\setlength{\tabcolsep}{3pt}
\begin{tabular}{@{}lccr@{}}
\toprule
\textbf{Shard} & $\Delta$F1 (pp) & $\Delta$EM (pp) & $\Delta$Tokens (\%) \\
\midrule
{[0:200]}     & $+0.44$ & $-3.50$ & $-36.9$ \\
{[200:400]}   & $+2.24$ & $-1.00$ & $-37.1$ \\
{[400:600]}   & $+2.28$ & $-3.00$ & $-37.3$ \\
\midrule
mean $\pm$ std & $+1.65 \pm 1.05$ & $-2.50 \pm 1.32$ & $-37.1 \pm 0.19$ \\
\bottomrule
\end{tabular}
\caption{Three-shard paired Stage-1 \texttt{fixed\_peer\_calibrated} vs.\ degenerate-Stage-2 \texttt{edo\_stage2\_chain} prototype on HotpotQA-200 with \texttt{gpt-4.1-mini}. Token cost reduction is consistent across shards ($\sigma{=}0.19$\%); F1 effect is inside noise; EM trends mildly negative.}
\label{tab:appendix-d-shard}
\end{table}

The token cost reduction is consistent across shards (cross-shard $\sigma{=}0.19$\%) and is the only effect that is statistically robust at this sample size; the F1 lift sits inside noise at $n{=}200$ per shard, and EM trends mildly negative ($-2.50 \pm 1.32$ pp). We report this as \textbf{preliminary and not statistically demonstrated}: the action-policy menu compresses hop count and token cost, but at this sample size the F1 trend is unresolved and the EM regression is a real (if mild) cost --- the prototype is not yet a Pareto improvement over Stage-1 across all metrics simultaneously. The headline empirical claims of the paper are deliberately not based on this preliminary evidence. The Stage-2 evaluation agenda described in \S 4.4 (E1--E5), executed at full validation scale ($n{=}7405$) with $\geq 3$ seeds and paired bootstrap CI, is required to settle whether the F1 trend is real and whether the EM regression reflects a missing synthesizer pass that future Stage-2 iterations must restore.

\begin{table}[h]
\centering
\small
\setlength{\tabcolsep}{3pt}
\begin{tabular}{@{}p{0.38\columnwidth}p{0.52\columnwidth}@{}}
\toprule
\textbf{Item} & \textbf{Safe status} \\
\midrule
HotpotQA-200 & Main evidence; single seed only. \\
Backbone sensitivity & Supported by chain-200 artifacts. \\
Cost-normalised F1 & Efficiency view only; not a raw-F1 win. \\
Memory/tool prototype & Infrastructure; no performance claim. \\
Error taxonomy export & Reproducible tooling; no manual taxonomy yet. \\
\bottomrule
\end{tabular}
\caption{Compact evidence-status checklist. The table separates supported claims from evidence still required for stronger EDO / Best-Paper claims.}
\label{tab:evidence-status}
\end{table}

\begin{figure}[h]
  \centering
  \includegraphics[width=0.88\columnwidth]{fig4_cost_normalized_f1.pdf}
  \caption{Cost-normalised F1 view for the current HotpotQA artifacts. This appendix figure preserves the efficiency-axis placeholder without consuming main-body pages.}
  \label{fig:cost-normalized-f1}
\end{figure}

\begin{figure}[h]
  \centering
  \includegraphics[width=0.88\columnwidth]{fig2_backbone_sensitivity.pdf}
  \caption{Backbone-sensitive Stage-1 ordering on the chain-200 HotpotQA slice. The ordering inverts under the stronger backbone, so the main text treats this as a scope guard rather than a finalized full-validation ranking.}
  \label{fig:backbone-sensitivity}
\end{figure}

\section*{Appendix E: EDO Stage-2 Execution Loop (Pseudocode)}
\label{sec:appendix-e}

This appendix gives the full pseudocode for the EDO Stage-2 execution loop referenced from \S 3.6, \S 3.9, and Finding~3 in \S 4.3. Section 3.6 of the main body specifies the prose semantics and the Stage-1 TCPB scoring formula; this appendix records the operational form. Empirical numbers in the main paper come from the degenerate Stage-1 instance only --- \textsc{Split} is disabled, \textsc{Audit} defaults to \textsc{Accept}, and the vector belief is collapsed via mean-axis fold (Prototype Scope Box, item (v)). Deterministic instantiations of $\phi$, $\textsc{Fit}$, and $\textsc{evidence\_extract}$ are given in the executable spec [Anonymous Suppl.] \S 1--3.

\begin{algorithm}[!ht]
\caption{EDO Stage-2 execution loop (single sample, recursive). The full EDO mechanism integrating R1 \texttt{split}, R2 recursive audit, and R3 vector-belief routing. The main paper reports the degenerate Stage-1 instance (TCPB), where \textsc{Split} is disabled and \textsc{Audit} defaults to \textsc{Accept}.}
\label{alg:edo-stage2-loop}
\small
\begin{algorithmic}[1]
\Require Root task $z_0$, gold $y$; agent set $\mathcal{A}$ with sparse adjacency $\mathcal{N}$; vector belief $B_i^{t}{:}\,\mathcal{A}{\times}\mathcal{A}{\to}[0,1]^{7}$ (R3); empty task tree $\mathcal{T}$ rooted at $z_0$; constants $H_{\max}{=}4$, $F_{\textsc{th}}{=}0.5$, $\delta_a{=}0.02$, $\nu{=}0.2$, $k_{\max}{=}3$, $d_{\max}{=}3$, $n_{\max}{=}12$.
\Ensure Predicted $\hat{y}$; updated belief $B^{t+1}$.
\Function{Process}{$z, i, \mathcal{T}$}
  \State $\phi \gets \textsc{extract}(z)$ \Comment{rule, no LLM call}
  \State $U^{\textsc{s}} \gets \textsc{Fit}(P_i,\phi) - \lambda_c\textit{Cost}_{\textit{self}}$
  \For{$j \in \mathcal{N}(i)$}
    \State $U^{\textsc{o}}_j \gets \textsc{Fit}(B_i^{t}(j),\phi) - \lambda_a\textit{AuditCost}$
  \EndFor
  \State $U^{\textsc{p}} \gets \textsc{SplitGain}(z) - \lambda_m\textit{MergeCost}$ \Comment{R1}
  \If{$\textit{depth}(z){\geq}d_{\max}$ \textbf{or} $|\mathcal{T}|{\geq}n_{\max}$}
    \State $U^{\textsc{p}} \gets {-}\infty$ \Comment{bound: SPLIT disabled}
  \EndIf
  \State $a \gets \arg\max\{U^{\textsc{s}},\,\max_j U^{\textsc{o}}_j,\,U^{\textsc{p}}\}$
  \If{$a{=}\textsc{DoSelf}$}
    \State \Return $\textsc{LLM\_answer}(z,P_i),\,\textsc{Accept}$
  \ElsIf{$a{=}\textsc{Outsource}(j^{\star})$}
    \State $r,\textit{stat} \gets \textsc{Process}(z, j^{\star}, \mathcal{T})$
    \State $e \gets \textsc{audit}(i,j^{\star},z,r)$ \Comment{R2; rule-based; no LLM call by default}
    \State $\mathcal{T}.\textsc{apply\_audit}(z,e)$ \Comment{writes \texttt{audit\_status}; jsonl-buffered}
    \If{$e.\textit{decision}{\in}\{\textsc{Accept},\textsc{AcceptWithNote}\}$} \Return $r, e.\textit{decision}$ \EndIf
    \If{$e.\textit{decision}{=}\textsc{RejectReroute}$ \textbf{and} unvisited $j' \in \mathcal{N}(i)$ exists}
      \State \Return \textsc{Process}$(z, j', \mathcal{T})$
    \Else \, $a \gets \textsc{Split}$ \Comment{escalate to RESPLIT}
    \EndIf
  \EndIf
  \If{$a{=}\textsc{Split}$}
    \State $\{z_1,\!\dots\!,z_k\} \gets \textsc{LLM\_decompose}(z)$ with $k{\leq}k_{\max}$ \Comment{R1}
    \State $\mathcal{T}.\textsc{add\_children}(z,\{z_l\})$
    \State \Return $\textsc{integrate}(\{\textsc{Process}(z_l,i,\mathcal{T})\}_{l=1}^k),\,\textsc{Accept}$
  \EndIf
\EndFunction
\State $\hat{y},\_ \gets \textsc{Process}(z_0, \textit{root\_agent}, \mathcal{T})$
\Statex \quad\textit{// R3 vector-belief update (one EMA step per audit event)}
\For{$e \in \mathcal{T}.\textit{audit\_events}$}
  \State $B_i^{t{+}1}(j) \gets \textsc{clip}\!\big((1{-}\nu) B_i^{t}(j) + \nu\cdot \textsc{evidence\_extract}(e,\phi(z)),0.05,0.95\big)$
\EndFor
\Statex \quad\textit{// Stage-1 fallback: terminal scalar update on root accepting agent (TCPB)}
\State $r \gets {+}1$ \textbf{if} $F_1(\hat{y},y){\geq}F_{\textsc{th}}$ \textbf{else} ${-}1$;\, project $B^{t{+}1} \to c$ via mean-axis fold (cf.\ Prototype Scope Box (v))
\State \Return $\hat{y},\,B^{t+1}$
\end{algorithmic}
\end{algorithm}

\section*{Appendix F: Theoretical Supplement (Formal Statements and Proofs)}
\label{sec:appendix-f}

This appendix records the formal statements and proofs of the four structural results referenced in \S 3 (1-locality invariant, complexity bound, persona-tag convergence, TCPB-as-degenerate-EDO reduction). The main body gives concise statements and defers technical detail here.

\paragraph{F.1 Decentralization invariant (1-locality).} A core design constraint of EDO is that every agent $a_i$ decision at time $t$ depends only on $i$'s own state $S_i^{t-1}$ and the visible neighbour set $\mathcal{N}(i)$ (plus messages received along incoming edges). Formally, we require the selection $\pi_i(z)$ in \S 3.3 to factor as $\pi_i(z)=f_i(S_i^{t-1}, \{S_j^{t-1}: j\in\mathcal{N}(i)\}, z)$ without access to $\{S_k^{t-1}\}_{k\notin\{i\}\cup\mathcal{N}(i)}$, to $\mathcal{G}$ as a whole, or to any global aggregator. This \emph{1-locality} invariant is preserved by the three-action policy (self / outsource-to-neighbour / split-locally) because none of the three actions reads or writes non-neighbour state. The Stage-1 TCPB prototype satisfies 1-locality trivially since the \texttt{chain} topology makes $\mathcal{N}(i)$ a single-element set $\{i+1\}$; Stage-2 graphs (sparse random / small-world / community-bridge) preserve it by construction: the execution loop of Algorithm~\ref{alg:edo-stage2-loop} reads $B_i^t(j)$ only for $j\in\mathcal{N}(i)$ (lines 4--6), writes $B_i^{t+1}(j)$ only for the same $j$ range (lines 32--34), and never reads global graph state $\mathcal{G}$.

\paragraph{F.2 Complexity bound on Process($z$).} With the Stage-2 recursion bounds $(H_{\max}, d_{\max}, n_{\max}, k_{\max}){=}(4,3,12,3)$ (Algorithm~\ref{alg:edo-stage2-loop}), the call count admits a closed-form bound. A single invocation of \textsc{Process}($z$) issues at most one \textsc{Decompose} call, at most $k_{\max}$ subcalls, and, along any root-to-leaf path in the delegation tree, at most $d_{\max}$ nested \textsc{Split}s and at most $H_{\max}$ \textsc{Outsource} steps, with an additional per-\textsc{Outsource} \textsc{Audit} call. Hence the total LLM-call count per input is bounded by $T_{\mathrm{LLM}}(z) \le (H_{\max}{+}1)\cdot(k_{\max}^{\,d_{\max}}) \le 5\cdot 27 = 135$ in the worst case, and by a looser union-bound $\le n_{\max}\cdot(H_{\max}{+}2) = 72$ when the total task-tree size is capped by $n_{\max}$. In the degenerate Stage-1 TCPB instance (Split disabled, Audit $\equiv$ Accept), both bounds collapse to $T_{\mathrm{LLM}}^{\mathrm{TCPB}}(z) \le H_{\max}{=}4$ LLM calls per input, matching the observed chain-200 per-sample call count. These bounds are decidable at scheduling time and govern the provider-side parallelism and quota pre-flight estimates (Appendix~B).

\paragraph{F.3 Convergence of the persona-tag update under stationarity.} Let $V_i^t{=}\eta_{\mathrm{loc}}\mathrm{LocalValue}_i^t{+}\eta_{\mathrm{trm}}\mathrm{TerminalValue}_i^t{-}\eta_{\mathrm{rew}}\mathrm{ReworkPenalty}_i^t$ denote the composite social-value signal produced for agent $i$ at time $t$. Under two standard assumptions --- (A1) \emph{environmental stationarity}: the task distribution and neighbour interaction protocol are time-invariant, so the signals $V_i^t$ are i.i.d.\ draws with $\mathbb{E}[V_i^t]{=}\bar V_i$ and $\mathrm{Var}(V_i^t){<}\infty$; and (A2) \emph{axis-wise boundedness}: $V_i^t$ and the pre-clip state $(1-\mu)P_i^t+\mu V_i^t$ remain within the clip range $[P_{\min}, P_{\max}]^{7}$ almost surely --- the update $P_i^{t+1}{=}(1{-}\mu)P_i^t{+}\mu V_i^t$ is a standard exponentially-weighted moving average (EMA) with fixed rate $\mu\in(0,1)$. Its sample mean satisfies
\[
P_i^{t} = (1-\mu)^{t} P_i^{0} + \mu\sum_{s=1}^{t}(1-\mu)^{t-s}V_i^{s},
\]
so $\lim_{t\to\infty}\mathbb{E}[P_i^t]{=}\bar V_i$ and the process converges in $L^{2}$ to $\bar V_i$ (cf.\ Robbins--Monro, applicable here because the step size $\mu$ is constant but the summation is normalised by the geometric weights $(1-\mu)^{t-s}$ whose total mass is $1-(1-\mu)^{t}\to 1$). The clip $\mathrm{clip}(\cdot)$ is non-expansive, so it does not break $L^{2}$ convergence under (A2). Under stationarity, therefore, the persona-tag state $P_i^t$ admits a \emph{unique fixed point} $P_i^\infty{=}\bar V_i$ interpretable as the time-averaged socially-observed value of agent $i$. Non-stationary settings (e.g.\ an adversarial neighbour or a drifting task distribution) violate (A1); in that regime $P_i^t$ no longer converges but tracks a sliding window of recent value with effective horizon $\approx 1/\mu$ steps. The delivered Stage-1 TCPB prototype runs under (A1) on the HotpotQA slice and with $\mu{=}0.5$; stability audits across the 200 samples showed no divergence, so the bounded-variation assumption (A2) is satisfied in practice.

\paragraph{F.4 Reduction (TCPB as a degenerate EDO instance).} Let $\mathcal{E}_{\mathrm{EDO}}$ denote the execution schema defined by Algorithm~\ref{alg:edo-stage2-loop}, parametrised by $(\pi, \textsc{Audit}, \mathrm{Upd})$ where $\pi$ is the action policy (\S 3.3), $\textsc{Audit}$ is the 4-class decision protocol (\S 3.4), and $\mathrm{Upd}$ is the vector-belief update (\S 3.5). The delivered TCPB schema $\mathcal{E}_{\mathrm{TCPB}}$ is obtained from $\mathcal{E}_{\mathrm{EDO}}$ by the following three restrictions: (R$_\textsc{Split}$) the \textsc{Split} branch is disabled by setting its gate weight in $\pi$ to $-\infty$; (R$_\textsc{Audit}$) \textsc{Audit} is replaced by the constant map $\textsc{Audit}(\cdot){\equiv}\textsc{Accept}$; and (R$_B$) the vector belief $B_i^t(j)\in[0,1]^{7}$ is projected onto the scalar competence $c[i]\in[0,1]$ via the mean-axis fold $c[i]{=}\tfrac{1}{7}\sum_{k}B_i^t(j)[k]$ and updated terminally. Under (R$_\textsc{Split}$ $\land$ R$_\textsc{Audit}$ $\land$ R$_B$), $\mathcal{E}_{\mathrm{EDO}}$ reduces to $\mathcal{E}_{\mathrm{TCPB}}$, and all main-paper empirical numbers are produced under $\mathcal{E}_{\mathrm{TCPB}}$. Any experimental claim in \S 4 about R1, R2, or R3 mechanism value would require relaxing at least one of (R$_\textsc{Split}$, R$_\textsc{Audit}$, R$_B$); none of those relaxations are delivered here (the main-paper Stage-2 roadmap).

\paragraph{F.5 R2 audit information-theoretic lower bound ($\Omega(H)$).}
We argue that per-hop recursive audit (R2) provides asymptotically more credit-assignment signal than terminal-only audit. Consider an $H$-hop delegation chain where each hop $h\in\{1,\ldots,H\}$ emits a local audit outcome $a_h\in\{0,1\}$ (accept / reject-class) and the terminal also emits a global outcome $y\in\{0,1\}$ (correct / incorrect). Let $f_h\in\{0,1\}$ be the latent ``true'' per-hop failure indicator. Assume each audit decision is \emph{$(1/2+\varepsilon)$-accurate}: $\Pr[a_h = f_h]\geq 1/2+\varepsilon$ for some $\varepsilon\in(0,1/2]$. Under R2, the learner observes the full audit sequence $\{a_1,\ldots,a_H, y\}$ per trial; under terminal-only audit (TCPB Stage-1 reduction, \S F.4), the learner observes $y$ only and must attribute failure to one of $H$ hops by uniform prior. Let $I(X;Y)$ denote mutual information. Then for credit assignment to the per-hop failure variable $F{=}(f_1,\ldots,f_H)$:
\[
I(F;\{a_1,\ldots,a_H,y\})\;\geq\; H\cdot\big(1 - H_2(\tfrac{1}{2}-\varepsilon)\big)\;\geq\;\Omega(H\varepsilon^2),
\]
where $H_2$ denotes binary entropy, while
\[
I(F;y)\;\leq\;H_2(\Pr[y{=}1])\;\leq\;1.
\]
Thus R2 recursive audit yields $\Omega(H)$ bits of credit signal per trial versus $O(1)$ for terminal-only audit, a gap that grows linearly in chain depth. This gives the $\Omega(H)$ sample-efficiency advantage referenced in \S 3.4 (``why recursive audit dominates terminal audit''). Proof sketch: the lower bound on $I(F;\{a_h\})$ follows from Fano's inequality applied per-hop; the upper bound on $I(F;y)$ follows from the data-processing inequality along the terminal channel. Full derivation in the anonymous supplement. The theorem is conditional on (i) audit $(1/2+\varepsilon)$-accuracy and (ii) per-hop failures are marginally independent; assumption (i) is empirically measurable via a separate audit-calibration protocol (Stage-2 experimental agenda E2, \S 4.4).

\paragraph{F.6 Formal metric family for H3 (specialisation entropy + persona divergence).}
To make H3 in the main-paper Stage-2 roadmap operational, let the non-negative 7-axis persona vector of agent $i$ at evaluation time be $P_i\in\mathbb{R}_{\ge 0}^{7}$, and define the normalised axis distribution
\[
q_i(k)\;=\;\frac{P_i(k)}{\sum_{\ell=1}^{7}P_i(\ell)},\qquad k\in\{1,\ldots,7\},
\]
with the convention $q_i(\cdot){=}(1/7,\ldots,1/7)$ if $\sum_{\ell}P_i(\ell){=}0$. The \emph{specialisation entropy} of agent $i$ is
\[
H_{\mathrm{spec}}(i)\;=\;-\sum_{k=1}^{7} q_i(k)\log q_i(k),
\]
and the society-level summary used by H3 is the mean $\bar H_{\mathrm{spec}}{=}\frac{1}{|\mathcal{A}|}\sum_i H_{\mathrm{spec}}(i)$. Lower values indicate sharper axis concentration (more specialised agents), while the uniform baseline gives $H_{\mathrm{spec}}(i){=}\log 7$. We pair this with a \emph{persona-divergence matrix}
\[
D_{ij}\;=\;1-\frac{\langle P_i,P_j\rangle}{\|P_i\|_2\|P_j\|_2},
\]
again treating zero vectors as divergence $0$. H3 is therefore evaluated not by one scalar alone but by the joint pattern: (i) $\Delta \bar H_{\mathrm{spec}}<0$ versus the scalar-TCPB baseline and (ii) a non-degenerate off-diagonal spread in $\{D_{ij}\}$ that exceeds a random axis-permutation null. The random-permutation baseline keeps the marginal axis mass of each $P_i$ fixed but permutes axis labels independently across agents; significance is then tested on $\Delta \bar H_{\mathrm{spec}}$ and on the mean pairwise divergence under that null.

\paragraph{F.7 Deterministic Stage-1 extractor for $\phi(z)$.}
The Stage-1 TCPB implementation does not issue an extra LLM call to instantiate $\phi(z)$. Instead, it uses a deterministic proxy from the current packet state: the decomposition coordinate turns on for root-level questions that trigger the multi-hop heuristic; the verification coordinate turns on once hop depth reaches $2$; the integration coordinate turns on once the visible evidence set reaches size $8$; the exploration and evidence-breadth coordinates are simple monotone functions of evidence count; uncertainty is copied from the packet; and cost sensitivity is the normalised hop depth $\min(1,\texttt{hop\_count}/H_{\max})$. The multi-hop heuristic itself is also deterministic: a root question is marked multi-hop if it contains one of a small set of bridge/comparison indicators (e.g.\ ``both'', ``and'', ``which of'') or if two or more capitalised proper-noun tokens appear after the first token. This rule set is the executable proxy behind $\phi(z)$ in \S 3.3 and matches the released routing-feature builder.
